% =====================================================================
%  Penalized vs. Constrained Wasserstein Adversarial Training
%  in Deep Hedging  --  REPORT SKELETON
%
%  * Grey boxes (\guide{...}) say what belongs in each part. Hide all of them
%    at once with \showguidesfalse below.
%  * Dashed boxes (\placeholderfig) mark figures to produce.
%  * Red [TODO] marks things to fill in or verify.
%  * Written equations are the settled formulation; check against the code.
% =====================================================================
\documentclass[11pt]{article}

\usepackage[margin=2.5cm]{geometry}
\usepackage{amsmath,amssymb,amsthm,bm}
\usepackage{booktabs}
\usepackage{graphicx}
\usepackage{xcolor}
\usepackage{algorithm}
\usepackage{algpseudocode}
\usepackage[colorlinks=true,linkcolor=blue!50!black,citecolor=blue!50!black,urlcolor=blue!50!black]{hyperref}
\usepackage[capitalise,noabbrev]{cleveref}

% ---------- guidance toggles ----------
\newif\ifshowguides
\showguidestrue            % <-- \showguidesfalse hides all guidance
\newcommand{\guide}[1]{\ifshowguides
  \par\smallskip\noindent\fcolorbox{gray!50}{gray!6}{\parbox{\dimexpr\linewidth-2\fboxsep-2\fboxrule}{\small\itshape\color{gray!80!black}#1}}\par\smallskip\fi}
\newcommand{\placeholderfig}[3]{% {caption}{label}{what to plot}
  \begin{figure}[htbp]\centering
  \fbox{\parbox[c][4.5cm][c]{0.8\linewidth}{\centering\small\itshape\color{gray!70!black}FIGURE TO PRODUCE:\\[3pt] #3}}
  \caption{#1}\label{#2}\end{figure}}
\newcommand{\todo}[1]{\textcolor{red!70!black}{[TODO: #1]}}

% ---------- notation ----------
\newtheorem{proposition}{Proposition}
\newcommand{\E}{\mathbb{E}}
\newcommand{\Ups}{\Upsilon}
\newcommand{\cvar}{\mathrm{CVaR}}
\newcommand{\PnL}{\mathrm{PnL}}

\title{Penalized versus Constrained Wasserstein Adversarial Training\\in Deep Hedging\\[4pt]\large QF4204 Project Report}
\author{Your Name}
\date{\today}

\begin{document}
\maketitle

% =====================================================================
\begin{abstract}
\guide{Write LAST, 150--200 words: (1) deep hedges are fragile to model misspecification;
(2) \cite{he2025} train against the worst case in a Wasserstein ball of fixed radius $\delta$ and
report sensitivity to $\delta$; (3) we replace the hard constraint by a Lagrangian penalty
$\lambda$ with the same transport cost; (4) headline findings in numbers: matched performance,
sensitivity of out-of-distribution risk to $\lambda$ vs.\ $\delta$, cross-robustness.}
\end{abstract}

% =====================================================================
\section{Introduction}\label{sec:intro}
\guide{About 1 page. (a) Deep hedging \cite{buehler2019} learns hedges from simulated paths and
inherits simulator error. (b) Robust deep hedging: parameter uncertainty \cite{lss2022};
Wasserstein adversarial training \cite{he2025}, the strongest recent result. (c) Their stated
limitation: performance depends on the radius $\delta$, which may differ across models.
(d) Penalized Wasserstein DRO \cite{sinha2018} is the Lagrangian counterpart of the constrained
problem; with a fixed penalty the effective radius adapts to the model during training.
(e) Research question and contributions.}

\paragraph{Research question.}
\todo{Does replacing the hard Wasserstein constraint in adversarial deep hedging with a Lagrangian
penalty, under the same transport cost, preserve hedging robustness, and does it reduce the
sensitivity of out-of-distribution performance to the robustness hyperparameter?}

\paragraph{Contributions.}\mbox{}
\guide{Keep checkable, one-to-one with \cref{sec:results}:
(1) a penalized adversary for deep hedging with exactly the geometry of \cite{he2025}'s attack,
so only constraint vs.\ penalty differs;
(2) a first-order calibration $\lambda(\delta)=\Ups/(2\delta)$ from \cite{bartl2021} that matches
the two formulations' strength;
(3) a controlled comparison of in- and out-of-distribution CVaR, sensitivity to the
hyperparameter, and cross-robustness. Claim ``to our knowledge'' novelty only after checking
\cite{he2025}'s related work.}

% =====================================================================
\section{Background and related work}\label{sec:background}
\subsection{Deep hedging}
\guide{\cite{buehler2019}: network hedge positions, convex risk measure of terminal P\&L,
simulated paths; handles frictions and incomplete markets.}

\subsection{Robust deep hedging}
\guide{\cite{lss2022}: robustness through the training DATA (parameter intervals).
\cite{he2025}: robustness through the OBJECTIVE (worst case in a Wasserstein ball, solved by a
projected attack); its first-order structure comes from \cite{bartl2021,bai2023}. Your work
changes how that objective's ball is enforced.}

\subsection{Constrained and penalized Wasserstein DRO}
\guide{Define the Wasserstein ball and the worst-case expectation. State the Lagrangian dual
linking the constrained and penalized problems \cite{blanchet2019,gao2023} (\cref{eq:duality}).
Penalized adversarial training for neural networks: \cite{sinha2018}, with a convergence
guarantee when the loss is smooth in the input and $\lambda$ exceeds its curvature. Sensitivity
of the worst case in $\delta$: \cite{bartl2021}, used here for calibration.}

% =====================================================================
\section{Problem setup}\label{sec:setup}
\guide{Match the authors' code exactly. Notation: $h_t$ hedge positions, $\delta$ radius,
$\lambda$ penalty, $\alpha$ CVaR level, $\gamma$ clean-loss weight (code: \texttt{--alpha}).}

\subsection{Market model}
Under the Heston model \cite{heston1993},
\begin{equation}\label{eq:heston}
dS_t = \sqrt{v_t}\,S_t\,dW^1_t,\qquad
dv_t = \kappa(b - v_t)\,dt + \sigma\sqrt{v_t}\,dW^2_t,\qquad
d\langle W^1,W^2\rangle_t = \rho\,dt ,
\end{equation}
simulated with a Broadie--Kaya-type scheme. Variance risk is hedged with a variance swap whose
price $S^{(2)}_t$ is a closed-form function of the variance path (\cite{he2025}, App.~B).

\begin{table}[htbp]\centering
\caption{Simulation parameters (\texttt{Heston\_generator.py}).}\label{tab:params}
\begin{tabular}{lll}\toprule
Parameter & Symbol & Value\\\midrule
Initial spot / strike & $S_0$, $K$ & 100\\
Initial variance & $v_0$ & 0.04\\
Mean reversion & $\kappa$ & 1\\
Long-run variance & $b$ & 0.04\\
Vol-of-vol & $\sigma$ & 2\\
Correlation & $\rho$ & $-0.7$\\
Drift & $\mu$ & 0\\
Maturity / steps & $T$, $n$ & $30/365$, 30\\
CVaR level & $\alpha$ & 0.5\\\bottomrule
\end{tabular}
\end{table}

\subsection{Hedging objective}
Positions $h_t=f^{\theta_t}_t(\log S_t, v_t)\in\mathbb{R}^2$ (stock, variance swap), one network per
step. Short a call, $Z=(S_T-K)^+$; loss $X=Z-\PnL_\theta$, $\PnL_\theta=\sum_{t<n}h_t\cdot(\bm S_{t+1}-\bm S_t)$.
With the Rockafellar--Uryasev form of CVaR \cite{rockafellar2000},
\begin{equation}\label{eq:loss}
\ell(\theta,w;x)=w+\frac{\big(Z(x)-\PnL_\theta(x)-w\big)^+}{1-\alpha},\qquad
\cvar_\alpha(X)=\min_w \E[\ell(\theta,w;x)],
\end{equation}
with path $x=(S,v)$ and threshold $w$ (\texttt{p0}) trained jointly with $\theta$.

\subsection{Transport cost}
Per path and coordinate $c\in\{S,V\}$, with $V$ measured in units scaled by 100 and $t=0$ fixed,
\begin{equation}\label{eq:cost}
c(x,x')=\sum_{c\in\{S,V\}}\ \max_{1\le t\le n}|x_{c,t}-x'_{c,t}|^2,
\qquad W_c(P,Q)=\inf_{\pi\in\Pi(P,Q)}\E_\pi[c(x,x')] .
\end{equation}
\guide{Stress that BOTH formulations use exactly this cost; the only difference is how it is
bounded.}

% =====================================================================
\section{Method}\label{sec:method}

\subsection{Constrained formulation (baseline)}
\begin{equation}\label{eq:constrained}
\mathcal{V}_\delta(\theta,w)=\sup_{Q:\,W_c(P,Q)\le\delta^2}\E_Q[\ell(\theta,w;x)],\qquad
\min_{\theta,w_c,w_a}\ \gamma\,\E_P[\ell(\theta,w_c;x)]+\mathcal{V}_\delta(\theta,w_a).
\end{equation}
\guide{Approximated in \cite{he2025} by a 20-iteration budget attack with one shared budget across
paths, projected onto the ball (step size $4\delta/20$).}

\subsection{Penalized formulation (ours)}
\begin{equation}\label{eq:penalized}
\mathcal{V}^\lambda(\theta,w)=\sup_Q\big\{\E_Q[\ell]-\lambda W_c(P,Q)\big\}
=\E_P\Big[\sup_{x'}\big\{\ell(\theta,w;x')-\lambda\,c(x,x')\big\}\Big],
\end{equation}
\begin{equation}\label{eq:pentrain}
\min_{\theta,w_c,w_a}\ \gamma\,\E_P[\ell(\theta,w_c;x)]+\mathcal{V}^\lambda(\theta,w_a).
\end{equation}
The inner problem in \eqref{eq:penalized} separates across paths. Since $\lambda c$ does not
depend on $\theta$, the gradient of \eqref{eq:pentrain} is that of $\ell$ at the per-path
maximisers $x^*(x)$ (Danskin), so training differs from \eqref{eq:constrained} only in the attack.

\subsection{Relation between the two}
Under standard conditions \cite{blanchet2019,gao2023},
\begin{equation}\label{eq:duality}
\mathcal{V}_\delta(\theta,w)=\inf_{\lambda\ge0}\big\{\lambda\delta^2+\mathcal{V}^\lambda(\theta,w)\big\}.
\end{equation}
\guide{Interpretation: the constrained problem is the penalized one with a multiplier that is
re-optimised for every $\theta$. Fixing $\lambda$ instead fixes the PRICE of transport, so the
realised radius changes as the network changes. This is the mechanism behind research question
(2): it may make performance less sensitive to the hyperparameter, or weaken the adversary.}

\subsection{Calibrating $\lambda$ to $\delta$}
For the cost \eqref{eq:cost}, linearising $\ell$ gives the per-path best response
$b_{n,c}=\|\nabla_{x_c}\ell_n\|_1/(2\lambda)$ (norms over $t\ge1$), hence a realised radius
\begin{equation}\label{eq:calib}
\Big(\E\textstyle\sum_c b_{n,c}^2\Big)^{1/2}=\frac{\Ups}{2\lambda},\qquad
\Ups=\Big(\E\textstyle\sum_c\|\nabla_{x_c}\ell\|_1^2\Big)^{1/2},
\qquad\Longrightarrow\qquad \lambda(\delta)=\frac{\Ups}{2\delta},
\end{equation}
where $\Ups$ is the sensitivity of \cite{bartl2021} for this cost
($\mathcal{V}_\delta=\E\ell+\delta\Ups+o(\delta)$).
\guide{Derivation in \cref{app:derivation}, including the consistency check that plugging
\eqref{eq:calib} into \eqref{eq:duality} recovers $\E\ell+\delta\Ups$. Calibrate on the clean
network the robust phase starts from; it is a first-order match, verified empirically in
\cref{sec:res-calib}.}

\subsection{The penalized attack}
\begin{algorithm}[htbp]
\caption{Penalized attack for one batch (same parameterisation as \cite{he2025}'s budget attack).}\label{alg:pen}
\begin{algorithmic}[1]
\State $g\gets\nabla_x\ell$ at the clean paths; set $g_{\cdot,0}=0$
\State $s\gets\mathrm{sign}(g)$, \ $b_{n,c}\gets\|g_{n,c}\|_1/(2\lambda)$ \Comment{exact optimum of linearised problem}
\For{$k=1,\dots,K$}
  \State $D_{n,c,t}\gets b_{n,c}s_{n,c,t}$; evaluate $J_n=\ell_n(x+D)-\lambda c_n(D)$; keep best $D$ per path
  \State $s\gets\mathrm{clip}\big(s+0.75\,\mathrm{sign}(\partial\ell/\partial s)+0.25(s-s_{\mathrm{old}}),-1,1\big)$, $s_{\cdot,0}=0$
  \State $b\gets\max\big\{(1-\kappa)\,b+\kappa\,(\partial\ell/\partial b)/(2\lambda m^2),\,0\big\}$, \ $m=\max_t|s|$
\EndFor
\State \Return best $x+D$ per path, and realised radius $(\mathrm{mean}_n\,c_n)^{1/2}$
\end{algorithmic}
\end{algorithm}
\guide{Convergence (\cref{app:convergence}): if the loss is $\mu$-strongly concave along the
budget, the budget update converges to the optimum iff $\lambda>-\mu/2$, at rate
$1-\kappa(1+\mu/2\lambda)$ per step; below the threshold the inner problem is unbounded. ReLU,
the CVaR kink and the sup-norm cost mean the smoothness assumption of \cite{sinha2018} fails,
so state that no guarantee is claimed and check stability empirically.}

% =====================================================================
\section{Experimental design}\label{sec:experiments}
\subsection{Data}
Training sets of size $N$ (state which), each run on $10^5/N$ disjoint partitions, capped at
\todo{number} partitions. Test: first \todo{$10^5$} of $10^6$ in-distribution paths.
Out-of-distribution: 100 configurations $\times$ 10\,000 paths, each of $\kappa,b,\sigma,\rho$
scaled by a uniform factor in $[0.9,1.1]$.

\subsection{Training arms}
\guide{(i) Clean, matched control: authors' adversarial script with $\delta=0$ (same learning
rate and schedule as the robust arms; the authors' clean script uses a different learning rate,
0.05 vs.\ 0.005 --- report this). (ii) Constrained, $\delta$ from \cref{tab:hyper}. (iii)
Penalized, $\lambda$ calibrated by \eqref{eq:calib}, same $\gamma$. Then sweeps of $\delta$ and
$\lambda$ at one $N$. Same epochs (300 clean + 400 robust), optimiser, partitions for all.}

\begin{table}[htbp]\centering
\caption{Radius and clean-loss weight per training size (SV-attack), \cite{he2025} Table~5a.}\label{tab:hyper}
\begin{tabular}{lccccc}\toprule
$N$ & 5\,000 & 10\,000 & 20\,000 & 50\,000 & 100\,000\\\midrule
$\delta$ & 0.5 & 1.0 & 0.1 & 0.03 & 0.005\\
$\gamma$ (code: \texttt{alpha}) & 1 & 10 & 1 & 0 & 0\\
$\lambda(\delta)$, \eqref{eq:calib} & \todo{} & \todo{} & \todo{} & \todo{} & \todo{}\\\bottomrule
\end{tabular}
\end{table}

\subsection{Metrics}
\guide{(i) In-distribution $\cvar_{0.5}$; (ii) out-of-distribution $\cvar_{0.5}$, mean and worst
over the 100 configurations; (iii) $\cvar_{0.5}$ under the constrained attack at $\delta$ and under
the penalized attack at $\lambda$ (every model faces both); (iv) realised radius of each attack;
(v) training time (context only). Mean $\pm$ std over partitions.}

% =====================================================================
\section{Results}\label{sec:results}
\subsection{Calibration check}\label{sec:res-calib}
\guide{Output of \texttt{calibrate\_lambda.py --verify}: for each $\delta$, the calibrated $\lambda$ and
the ratio realised radius$/\delta$. Near 1 at small $\delta$ confirms \eqref{eq:calib}; the drift at
large $\delta$ measures nonlinearity.}

\subsection{Matched comparison}
\begin{table}[htbp]\centering
\caption{Hedging risk ($\cvar_{0.5}$, lower is better), mean $\pm$ std over partitions.}\label{tab:main}
\begin{tabular}{llcccc}\toprule
$N$ & Method & In-dist. & OOD mean & OOD worst & Attacked ($\delta$ / $\lambda$)\\\midrule
10k & Clean & & & & \\
    & Constrained & & & & \\
    & Penalized & & & & \\\midrule
20k & Clean & & & & \\
    & Constrained & & & & \\
    & Penalized & & & & \\\bottomrule
\end{tabular}
\end{table}

\subsection{Sensitivity to the robustness hyperparameter}
\placeholderfig{Out-of-distribution risk against realised radius, both formulations.}{fig:sensitivity}{x: realised radius (log scale), from sweeping $\delta$ (constrained) and $\lambda$ (penalized). y: OOD CVaR, mean over partitions with bands. Horizontal line: clean. Headline figure.}
\guide{Quantify ``flatness'', e.g.\ the spread of OOD CVaR over each grid, or the regret of a
mis-chosen hyperparameter relative to the best one. State the grids.}

\subsection{Cross-robustness}
\begin{table}[htbp]\centering
\caption{$\cvar_{0.5}$ of each trained model under each adversary.}\label{tab:cross}
\begin{tabular}{lcc}\toprule
Trained with $\downarrow$ / attacked by $\rightarrow$ & Constrained ($\delta$) & Penalized ($\lambda$)\\\midrule
Clean & & \\
Constrained & & \\
Penalized & & \\\bottomrule
\end{tabular}
\end{table}

\subsection{Training dynamics}
\placeholderfig{Realised radius of the adversary during the robust phase.}{fig:radius}{x: epoch 300--700. y: realised radius. Penalized at a few $\lambda$ (logged by \texttt{Heston\_train\_pen.py}); constrained is flat at $\delta$. Explains the mechanism behind \cref{fig:sensitivity}.}

\subsection{Computational cost (context)}
\guide{Seconds per robust-phase epoch for both. One sentence; not a research question.}

% =====================================================================
\section{Discussion}\label{sec:discussion}
\guide{When does the adaptive radius help or hurt? Does penalized training trade worst-case
protection for typical-case performance? What does cross-robustness say about the two notions
of robustness? Relate to \cite{he2025}'s limitation on radius sensitivity without overclaiming.}

\section{Limitations and future work}\label{sec:limitations}
\guide{(1) No convergence guarantee for the inner problem (non-smooth loss and cost); iteration
count fixed at 20; small $\lambda$ can be unbounded. (2) Calibration is first order. (3) Choice of
CVaR threshold inside the attack (batch quantile vs.\ trained $w_a$). (4) Simulated Heston only;
Black--Scholes not run. (5) One transport cost; a squared-$\ell_2$ cost would be a genuine geometry
change. (6) Plain Wasserstein on paths allows anticipative perturbations; causal transport is
the principled fix \todo{verify and cite causal-DRO sensitivity work}.}

\section{Conclusion}\label{sec:conclusion}
\guide{One paragraph with the headline numbers.}

% =====================================================================
\appendix
\section{Linearised best response and calibration}\label{app:derivation}
\guide{With $g=\nabla_x\ell$ and $D_{n,c}=b_{n,c}s_{n,c}$, H\"older gives
$\langle g_{n,c},D_{n,c}\rangle\le\|g_{n,c}\|_1 b_{n,c}$, attained at $s=\mathrm{sign}(g)$. Maximise
$\|g_{n,c}\|_1b-\lambda b^2$: $b^*=\|g_{n,c}\|_1/(2\lambda)$, value $\|g_{n,c}\|_1^2/(4\lambda)$. So
$\mathcal{V}^\lambda\approx\E\ell+\Ups^2/(4\lambda)$ and the realised radius is $\Ups/(2\lambda)$.
Consistency: in \eqref{eq:duality}, $\min_\lambda\{\lambda\delta^2+\Ups^2/(4\lambda)\}=\delta\Ups$ at
$\lambda=\Ups/(2\delta)$, recovering \cite{bartl2021}'s $\E\ell+\delta\Ups$. Include the numerical
verification (linear loss: objective $\Ups^2/(4\lambda)$, radius $\Ups/(2\lambda)$, random perturbations
never better).}

\section{Convergence of the budget update}\label{app:convergence}
\guide{Along a fixed sign profile let $\ell(b)=Gb-\tfrac{\mu}{2}b^2$. The update
$b\leftarrow(1-\kappa)b+\kappa(G-\mu b)/(2\lambda)$ has fixed point $G/(2\lambda+\mu)$ and factor
$1-\kappa(1+\mu/2\lambda)$: it converges iff $0<\kappa(1+\mu/2\lambda)<2$, i.e.\ (for $\kappa\le1$) iff
$\lambda>-\mu/2$. Below the threshold the objective is unbounded. Report the NumPy check.}

\section{Implementation and reproducibility}\label{app:impl}
\guide{Repository and commit hash; files added (\texttt{penalized.py}, \texttt{Heston\_train\_pen.py},
\texttt{calibrate\_lambda.py}, \texttt{evaluate.py}, \texttt{sensitivity.py}); authors' files unmodified;
attack in eval mode; threshold option; \texttt{--b\_max}; exact commands per table and figure; hardware
and run times. Note: the authors' attack always runs 20 iterations (its \texttt{iter} argument
only sets the step size).}

% =====================================================================
\begin{thebibliography}{99}

\bibitem{buehler2019}
H.~Buehler, L.~Gonon, J.~Teichmann, and B.~Wood.
\newblock Deep hedging.
\newblock \emph{Quantitative Finance}, 19(8):1271--1291, 2019.

\bibitem{he2025}
G.~He, T.~Sutter, and L.~Gonon.
\newblock Distributional adversarial attacks and training in deep hedging.
\newblock In \emph{Advances in Neural Information Processing Systems (NeurIPS)}, 2025.
\newblock arXiv:2508.14757.

\bibitem{lss2022}
E.~L{\"u}tkebohmert, T.~Schmidt, and J.~Sester.
\newblock Robust deep hedging.
\newblock \emph{Quantitative Finance}, 22(8):1465--1480, 2022.

\bibitem{sinha2018}
A.~Sinha, H.~Namkoong, and J.~Duchi.
\newblock Certifying some distributional robustness with principled adversarial training.
\newblock In \emph{International Conference on Learning Representations (ICLR)}, 2018.

\bibitem{blanchet2019}
J.~Blanchet and K.~Murthy.
\newblock Quantifying distributional model risk via optimal transport.
\newblock \emph{Mathematics of Operations Research}, 44(2):565--600, 2019.

\bibitem{gao2023}
R.~Gao and A.~Kleywegt.
\newblock Distributionally robust stochastic optimization with {W}asserstein distance.
\newblock \emph{Mathematics of Operations Research}, 48(2):603--655, 2023.

\bibitem{bartl2021}
D.~Bartl, S.~Drapeau, J.~Ob{\l}{\'o}j, and J.~Wiesel.
\newblock Sensitivity analysis of {W}asserstein distributionally robust optimization problems.
\newblock \emph{Proceedings of the Royal Society A}, 477:20210176, 2021.

\bibitem{bai2023}
X.~Bai, G.~He, Y.~Jiang, and J.~Ob{\l}{\'o}j.
\newblock Wasserstein distributional robustness of neural networks.
\newblock In \emph{Advances in Neural Information Processing Systems (NeurIPS)}, 2023.

\bibitem{rockafellar2000}
R.~T. Rockafellar and S.~Uryasev.
\newblock Optimization of conditional value-at-risk.
\newblock \emph{Journal of Risk}, 2(3):21--41, 2000.

\bibitem{heston1993}
S.~L. Heston.
\newblock A closed-form solution for options with stochastic volatility with applications to bond and currency options.
\newblock \emph{Review of Financial Studies}, 6(2):327--343, 1993.

\end{thebibliography}
\end{document}
